\chapter{Architecture and Methodology}\label{ch:architecture}
% Content cycle: C3. Intent: walk the compilation pipeline end to end; defend
% the two boundaries; explain the soundness gate and why it is sound here.

\section{The compilation pipeline}\label{sec:arch-pipeline}
% parse -> typed algorithm IR -> link -> ACFG -> transforms -> transfer/halo
% inference -> Petri-net IR -> soundness gate -> per-worker projection ->
% presentation-layer codegen. (Central figure.)

\section{The two boundaries}\label{sec:arch-boundaries}
% Algorithm/schedule; middle-end/presentation.

\section{Algorithm IR and linking}\label{sec:arch-link}

\section{The ACFG and its transforms}\label{sec:arch-acfg}
% vectorise, block, pipeline, reuse.

\section{Transfer and halo inference}\label{sec:arch-transfer}
% Automatic transfer synthesis, not automatic parallelism extraction.

\section{The Petri-net IR and the event contract}\label{sec:arch-petri}

\section{The soundness gate}\label{sec:arch-gate}
% Boundedness, deadlock, conflict-freedom by exact replay over the
% deterministic firing order; why this is sound for the restricted class.

\section{Per-worker projection}\label{sec:arch-projection}
% worker -> ordered EventList.

\section{Presentation-layer codegen}\label{sec:arch-codegen}

\todo{Write Chapter 6 (cycle C3).}
